\newcommand{\ModuleID}{EL-II.6}
\newcommand{\ModuleTitle}{Purpose, Result, and Indirect Command}
\newcommand{\ModuleStage}{Stage II — Core Grammar}
\newcommand{\ModuleUnit}{II.6}
\newcommand{\ModuleOutcome}{Classify outwardly similar \Latin{ut}-clauses from governor, signals, negation, sequence, and discourse force}
\newcommand{\ModulePrerequisites}{EL-II.5 mastery: identify subjunctive morphology and sequence before assigning use}
\newcommand{\ModuleSessionOne}{Build the purpose/result/indirect-command evidence matrix, including \Latin{ut/ne} and \Latin{ut non}, then complete Worksheet II.21 as the guided application.}
\newcommand{\ModuleSessionTwo}{Classify contextual clauses and transform positive and negative members of each family with preserved sequence, then complete Worksheet II.24 independently.}
\newcommand{\ModulePrintedPractice}{Worksheets II.21 (guided Variant A) and II.24 (independent Variant D), with terminal keys}
\newcommand{\ModuleExtensionPractice}{Worksheets II.22 and II.23 remain optional remediation or delayed-recall variants}
\newcommand{\ModuleNext}{EL-II.7, Substantive Clauses}
\newcommand{\ModuleMastery}{Classify at least eleven of twelve mixed clauses with a stated reason, select the correct negative marker in all eight test items, and compose one positive and one negative example of each family with correct sequence.}
\newcommand{\ModuleDelayNote}{}
\newcommand{\ModuleReferenceFocus}{%
  \begin{itemize}
    \item purpose: intended end, \Latin{ut/ne}, and relative-purpose alternatives;
    \item result: actual or represented outcome, result signals, and \Latin{ut/ut non};
    \item indirect command: a commanding, asking, warning, or persuading governor plus \Latin{ut/ne}; and
    \item sequence of tenses after primary and historic governors.
  \end{itemize}}
\newcommand{\ModuleContrast}{Negative purpose and indirect command ordinarily use \Latin{ne}; negative result uses \Latin{ut non}. Because positive clauses can all begin with \Latin{ut}, the governor, internal signals, and discourse relation must decide the family.}
\newcommand{\ModuleLessonSource}{\CourseRoot/shared/content/volume2/all}
\newcommand{\ModuleExerciseSource}{\CourseRoot/shared/exercises/volume2/all}
\newcommand{\ModuleWorkedBank}{II}
\newcommand{\ModuleExerciseBank}{II}
\newcommand{\ModuleWorksheetVariants}{A,D}
\newcommand{\ModuleScopeNote}{This packet teaches three central clause families without claiming that a connective alone supplies a complete parse.}
\newcommand{\ModuleEnrichment}{%
  \SubmoduleExtension
    {Three clause patterns}
    {An evidence matrix prevents \Latin{ut} from being treated as a one-word clause label.}
    {Create purpose, result, and indirect-command columns with rows for governor, discourse question, common signal, positive marker, negative marker, and sequence. Fill the matrix from the lesson before checking.}
    {Recite the negative marker for each family and one additional deciding feature.}
    {Purpose and indirect command use \Latin{ne} when negative; result uses \Latin{ut non}. Purpose answers intended end, result presents an outcome and often has a degree/result signal, and indirect command depends on a volitional or communicative governor.}
  \SubmoduleExtension
    {Reading the forms in context}
    {Classification strength increases as independent evidence converges.}
    {For each mixed lesson clause, assign one point each for a diagnostic governor, internal signal, expected negation, sequence, and discourse fit. Choose the narrowest family supported and mark any unresolved tie.}
    {Explain why \Latin{ut} plus subjunctive earns only partial evidence.}
    {The connective and mood establish a broad dependent pattern but are shared by multiple families. A defensible answer cites the governor and semantic relation, uses signals and negation as corroboration, and preserves ambiguity when the text does not distinguish purpose from result securely.}
  \SubmoduleExtension
    {Writing laboratory}
    {Positive-to-negative transformations test the clause family because the negative marker must change by family.}
    {Perform the lesson's three transformations. On each positive/negative pair, underline the governor, box the marker, and label sequence; do not change the intended relation.}
    {State which features remain constant when only polarity changes.}
    {The governor, clause family, relative time, and sequence remain constant; polarity and its marker change. Purpose and command pairs require \Latin{ne}; result requires \Latin{ut non}; all finite forms must retain correct sequence.}
}
